% gwebmac.tex -- plain-TeX macros for documents produced by gweave.
% This is the GWEB analogue of CWEB's cwebmac.tex.  Process a woven file with
% a plain-TeX engine, e.g.  pdftex foo.tex  (gwebmac.tex must be findable).

% ---- fonts ------------------------------------------------------------------
\font\titlefont=cmbx10 scaled\magstep2   % starred-section / page titles
\font\eightrm=cmr8                       % small roman for cross-reference notes
\font\eightit=cmti8

% ---- page layout ------------------------------------------------------------
\hsize=6.5in
\vsize=8.9in
\parindent=0pt
\parskip=0pt
\hbadness=10000 \hfuzz=300pt             % code lines may run wide; do not fuss
\tolerance=2000
\newdimen\Gindentunit \Gindentunit=1.4em % one indentation step in displayed code
\newdimen\Gtocindent  \Gtocindent=1.5em  % one indentation step in the contents
\newskip\Gsecskip     \Gsecskip=11pt plus 4pt minus 2pt
% (plain TeX already provides \S for the section sign used in the notes below.)

% A starred section records itself in \jobname.toc as it is shipped out, so that
% \Gcon can build a table of contents in a single TeX pass.  The page number is
% captured by the deferred \write; the title is stashed in a control sequence so
% it need not survive being written to (and read back from) a file.
\newwrite\Gtoc
\immediate\openout\Gtoc=\jobname.toc

% ---- section openers --------------------------------------------------------
% \GM{n} : an ordinary (unstarred) section, run-in number.
\def\GM#1{\par\penalty-50\vskip\Gsecskip\noindent
  \hbox{\bf #1.\enspace}\ignorespaces}

% \GN{depth}{n}{title} : a starred section.  A depth-0 section begins a new page,
% except for the very first one, which follows any title banner in the limbo.
\newif\ifGfirst \Gfirsttrue
\def\GN#1#2#3{\par
  \ifnum#1=0
    \ifGfirst \Gfirstfalse \vskip 6pt \else \penalty-300\vfil\eject \fi
  \else \penalty-150\vskip 14pt plus 4pt \fi
  % Stash the title and record a contents line *after* any page break above, so
  % the deferred \write captures the page this heading actually lands on.
  \expandafter\gdef\csname Gtoctitle\number#2\endcsname{#3}%
  \write\Gtoc{\string\GTline{#1}{#2}{\the\pageno}}%
  \noindent{\titlefont #2.\enspace #3.}\par\nobreak\vskip 4pt
  \noindent\ignorespaces}

% ---- displayed Go code ------------------------------------------------------
% \GB ... \GE bracket a code display; \GL typesets one line of it.
\def\GB{\par\nobreak\vskip 3pt\begingroup
  \baselineskip=11pt \rightskip=0pt plus 1fil \parfillskip=0pt plus 1fil}
\def\GE{\par\endgroup\vskip 3pt}
\def\GL#1#2{\par\noindent\kern#1\Gindentunit$#2$}

% ---- code token fonts (called inside the math of \GL) -----------------------
\def\GID#1{\hbox{\it #1\/}}      % identifier
\def\GKW#1{\hbox{\bf #1}}        % reserved word or predeclared type
\def\GST#1{\hbox{\tt #1}}        % string, rune, or verbatim text
\def\GCM#1{\hbox{\rm #1}}        % comment
\def\GNU#1{\hbox{\rm #1}}        % numeric literal

% ---- named-section references and headlines ---------------------------------
\def\Gangle#1#2{$\langle\,$#1\penalty100\thinspace{\sevenrm #2}$\,\rangle$}
\def\GX#1#2{\hbox{\Gangle{#2}{#1}}}        % inline reference in code
\def\GD#1#2{\par\vskip 3pt\noindent\hbox{\Gangle{#2}{#1}${}\equiv{}$}\hskip 3pt}
\def\GDp#1#2{\par\vskip 3pt\noindent
  \hbox{\Gangle{#2}{#1}${}\mathrel{+}\equiv{}$}\hskip 3pt}

% ---- cross-reference notes printed under a definition -----------------------
\def\GU#1{\par\nobreak\vskip 2pt\noindent{\eightrm This code is used in \S#1.}\par}
\def\GA#1{\par\nobreak\vskip 2pt\noindent{\eightrm This code is also defined in \S#1.}\par}
\def\GUL#1{$\underline{\hbox{#1}}$} % an underlined (defining) section number

% ---- index ------------------------------------------------------------------
\def\Ginx{\par\vfil\eject\centerline{\titlefont Index}\vskip 14pt
  \begingroup\parindent=0pt\parskip=2pt plus 1pt\rightskip=0pt plus 2em\relax}
\def\GII#1#2{\noindent#1:\quad #2.\par}      % one index entry: head and sections
\def\GIR#1{{\rm #1}}                          % @^...@>  roman index entry
\def\GIT#1{{\tt #1}}                          % @....@>  typewriter index entry
\def\9#1{{\rm #1}}                            % default for @:...@> entries
\def\GIC#1{\9{#1}}

% ---- list of named sections (refinements) -----------------------------------
\def\Gfin{\par\endgroup\vfil\eject
  \centerline{\titlefont List of Refinements}\vskip 14pt
  \begingroup\parindent=0pt\parskip=3pt plus 1pt\rightskip=0pt plus 2em\relax}
\def\GNS#1#2#3{\noindent\hbox{$\langle\,$#1$\,\rangle$}\quad
  {\eightrm defined in \S#2; used in \S#3.}\par}

% ---- table of contents and closing -----------------------------------------
% \GTline{depth}{n}{page} typesets one contents row; the title comes from the
% control sequence stashed by \GN.  It is defined locally inside \Gcon.
\def\Gcon{\par\endgroup            % close the refinements group opened by \Gfin
  \vfil\eject                      % ship the last page so all \write's execute
  \immediate\closeout\Gtoc
  \centerline{\titlefont Table of Contents}\vskip 14pt
  \begingroup
    \parindent=0pt \parfillskip=0pt \rightskip=0pt \parskip=2pt plus 1pt
    \def\GTline##1##2##3{\par\noindent\kern##1\Gtocindent
      \hbox{\bf ##2.}\enspace\csname Gtoctitle\number##2\endcsname
      \enspace\dotfill\enspace\hbox{##3}\par}%
    \input \jobname.toc
  \endgroup}
